% ------------------------------------------------------------
% CHAPTER 30
% ------------------------------------------------------------

\chapter{Derived Geometry and Discrete Measurement}

The analysis developed in previous chapters has revealed quantization as a form of symbolic collapse, delta–sigma feedback as a lamphrodynamic entropy-routing engine, and the modulator as an active-inference system that maintains coherence under continual perturbation. These interpretations converge naturally on the mathematics of derived geometry, where discrete observations arise from homotopy-corrected intersections, and measurement is treated not as a simple evaluation map but as a derived fiber product between continuous and discrete structures. In this chapter we develop the formal connection between derived geometry and delta–sigma modulation, showing that the quantizer, the loop filter, and the error dynamics admit a precise analog in the language of derived stacks, cotangent complexes, and higher homotopies.

At the heart of derived geometry lies the recognition that intersections between spaces often contain hidden structure. A classical intersection may appear transverse or degenerate, but the homotopical data lurking beneath the classical picture captures the infinitesimal relationships among the spaces. The derived fiber product encodes these relationships through chain complexes and higher-order compatibility conditions. In the context of measurement, a continuous field \(\Phi(t)\) and a discrete symbol space \(\mathcal{A}\) intersect along the quantizer map
\[
Q : M \longrightarrow \mathcal{A},
\]
where \(M\) is the underlying manifold of signal values or internal states. The quantizer selects a symbol \(a \in \mathcal{A}\) such that the continuous value lies within the collapse region associated with \(a\). In classical geometry, this would correspond to choosing a discrete equivalence class. In derived geometry, however, the collapse is understood as a derived intersection
\[
M \times^{\mathbf{R}}_{\mathcal{A}} \{a\},
\]
which encodes far more than the mere fact of collapse. It carries with it the infinitesimal relationships between \(M\) and \(\mathcal{A}\), including the constraint structure, the tangent and cotangent spaces, and the higher-order failure of transversality.

This structure corresponds precisely to what the delta–sigma modulator reconstructs through feedback. The classical quantizer discards information by collapsing \(x\) to \(Q(x)\). The derived viewpoint reveals that the discarded information is not lost but transformed into a higher-order complex that the modulator’s feedback loop attempts to recover. The cotangent complex of the derived fiber product captures the infinitesimal deviation between the continuous trajectory and the symbolic representation, and the loop filter acts as a homotopy lift that reconstructs the continuous path from the symbolic collapse.

To see this explicitly, let the cotangent complex of the derived fiber product be denoted by
\[
\mathbb{L}_{M \times^{\mathbf{R}}_{\mathcal{A}} \{a\}},
\]
which encodes the space of allowable infinitesimal variations of the collapse. The quantization error \(x[n] - Q(x[n])\) corresponds to a representative of a class in this cotangent complex. Linearization of the feedback dynamics produces an update equation
\[
e[n+1] = A e[n] + \delta[n],
\]
where \(\delta[n]\) is the new error injected by the quantization event. Interpreted through the derived lens, this recurrence computes a differential on the cotangent complex. The operator \(A\) corresponds to the action of the loop filter on infinitesimal deviations, while \(\delta[n]\) corresponds to the obstruction class representing the failure of the classical collapse to be smooth.

Each quantization event therefore produces an obstruction class in the derived intersection, and each iteration of the loop filter computes a homotopy that reduces this obstruction. The loop is a homotopy-correcting mechanism: it enforces coherence by lifting the symbolic collapse back into the continuous manifold, correcting for the loss of structure that occurs at each sampling instant. In the engineering language, this is noise shaping; in the geometric language, it is derived homotopy lifting.

The connection becomes even clearer when we interpret the delta–sigma loop as a derived mapping stack. A mapping stack between two spaces parametrizes maps modulo homotopy. In the case of delta–sigma modulation, the mapping from the continuous field \(\Phi(t)\) to the symbolic stream \(y[n]\) is not simply a function but a homotopy-coherent structure. The discrete symbol sequence represents a homotopy class of approximations to \(\Phi\). The delta–sigma loop filter refines this approximation by adjusting the internal state so that the mapping becomes increasingly accurate in a derived sense. One may therefore express the modulator as a sequence of homotopy lifts in the mapping stack
\[
\mathrm{Map}(M, \mathcal{A}),
\]
with each lift corresponding to a feedback iteration.

Derived geometry also clarifies why delta–sigma modulators achieve such high accuracy despite their brutally discontinuous quantizers. Classical quantization discards high-order information irreversibly. Derived collapse retains this information in the form of hidden homotopies. These homotopies can be recovered by applying smoothing operators that propagate the hidden structure forward in time. The loop filter acts precisely in this role: it is the operator that lifts the derived collapse back toward the continuous field. The classical quantizer does not annihilate structure; it merely relocates it into higher complexes, which the modulator gradually reintegrates into its internal state.

This viewpoint also sheds light on overload and nonlinear instability. When the signal moves outside the collapse region associated with the current symbol, the derived fiber product becomes highly singular. The cotangent complex grows in dimension or develops nontrivial higher homotopies. The feedback loop may fail to compute the necessary homotopy lift, resulting in error accumulation, runaway behavior, or limit cycles. These instabilities appear as derived obstructions: the failure of certain higher-order compatibility conditions that must hold for the homotopy lift to exist. This geometric interpretation explains why high-order modulators must be carefully designed to avoid pathological loop dynamics: their derived stacks contain richer obstruction theories with more stringent coherence conditions.

The derived perspective offers a path toward generalization. One may construct quantizers whose symbol spaces \(\mathcal{A}\) have nontrivial topological or algebraic structure, such as trees, hyperbolic spaces, or semantic manifolds. The collapse map then becomes a morphism between structured spaces, and the feedback loop becomes a derived map that preserves coherence across these structures. This opens the possibility of semantic quantization in which symbols encode higher meaning rather than simply numerical value. The derived mapping stack then becomes the natural domain for designing, analyzing, and generalizing such quantizers.

In summary, derived geometry reveals the hidden structure behind delta–sigma quantization: collapse is a derived fiber product, error is an element of the cotangent complex, feedback is a homotopy lift, stability is the absence of obstructions, and noise shaping is the iterative resolution of derived incompatibilities. The modulator becomes a homotopy-corrected measurement engine that preserves continuity through discrete collapse. This formalism provides a foundation for understanding measurement not only in engineering but across physics and cognition, where derived structures govern the interaction between continuous fields and discrete symbolic systems.

The next chapter develops the thermodynamic implications of this perspective, treating delta–sigma modulators as entropy engines whose operation reflects the universal laws governing energy, entropy, and information flow.


